\documentclass[compacto,10pt]{pubciencias-notas}

% -- Paquetes adicionales
\usepackage{comandosEPN}
\usepackage{array,xparse}
\usepackage{amssymb}
% -- Datos del libro
\autor{Departamento de Formación Básica}
\materia{Álgebra Lineal}
\nota{Semana 03 Clase 01}
\fecha{Semestre 2025-A}

\logouno[1.1cm]{Logo_EPN.eps}
\logodos[1.7cm]{Logo_DFB.eps}

\definecolor{colordef}{cmyk}{0.81,0.62,0.00,0.22}

\begin{document}

\encabezado

\vspace*{-8mm}

\setcounter{section}{1}
\section{Sistemas de ecuaciones}

\begin{defi}[Sistema de ecuaciones]
    Un sistema de ecuaciones lineales es un conjunto de $m$ ecuaciones lineales en las cuales figuran $n$ incógnitas a las que notaremos por: $x_1, x_2, \ldots, x_n$. Así, dados $a_{ij}\in\R$ y $b_i\in \R$ con $i\in I$ y $j\in J$, fijos, un sistema de ecuaciones es:
    \[
    \l\{
    \begin{array}{cccccccccc}
        a_{11}x_1& + &a_{12}x_2& + &\cdots &+& a_{1n}x_n &= &b_1,\\
        a_{21}x_1& + &a_{22}x_2& + &\cdots &+& a_{2n}x_n &= &b_2,\\
                &&&    &            &        & &\hspace{2mm} \vdots&\\
        a_{m1}x_1& + &a_{m2}x_2& + &\cdots &+& a_{mn}x_n &= &b_m.
    \end{array}
    \r.
    \]
\end{defi}

\begin{ejem}
El siguiente sistema de ecuaciones:
\[
    \systeme[x_1 x_2 x_3 x_4]{x_1 +x_2 +5x_3 -4x_4=3, x_1 -2x_2 +x_3 -2x_4=1, 3x_1 +x_2 -2x_3 -4x_4=-3.}
\]
Es un sistema de ecuaciones lineales de 3 ecuaciones con 4 incógnitas $x_1, x_2, x_3$ y $x_4$. 
\end{ejem}
%%%%%%%%%%%%%%%%%%%%%%%%%%%%%%%%%%%%%%
\subsection{Representación matricial de un sistema de ecuaciones}
%%%%%%%%%%%%%%%%%%%%%%%%%%%%%%%%%%%%%%
\begin{defi}[Representación matricial]
Dado un sistema de ecuaciones lineales de $m$ ecuaciones lineales en las cuales figuran $n$ incógnitas:
\[
\l\{
\begin{array}{cccccccccc}
    a_{11}x_1& + &a_{12}x_2& + &\cdots &+& a_{1n}x_n &= &b_1,\\
    a_{21}x_1& + &a_{22}x_2& + &\cdots &+& a_{2n}x_n &= &b_2,\\
            &&&    &            &        & &\hspace{2mm} \vdots&\\
    a_{m1}x_1& + &a_{m2}x_2& + &\cdots &+& a_{mn}x_n &= &b_m.
\end{array}
\r.
\]
donde $a_{ij}\in\R$ y $b_i\in \R$ con $i\in I$ y $j\in J$. A la matriz
\[
    A=\begin{pmatrix}
    a_{11} & a_{12} & \cdots & a_{1n}\\
    a_{21} & a_{22} & \cdots & a_{2n}\\
    \vdots & \vdots & \ddots & \vdots\\
    a_{m1} & a_{m2} & \cdots & a_{mn}
    \end{pmatrix}
\]
se la llama matriz de coeficientes del sistema y a 
\[
    b=\begin{pmatrix}
    b_1 \\ b_2 \\ \vdots \\ b_m
    \end{pmatrix}
    \texty
    x=\begin{pmatrix}
    x_1 \\ x_2 \\ \vdots \\ x_n
    \end{pmatrix}
\]
se las llama columnas de constantes y de incógnitas, respectivamente. Además, se dice que
\[
    Ax=b
\]
es la representación matricial del sistema de ecuaciones.
\end{defi}

\begin{ejem}\hspace{0pt}\label{Sem2:EjeSisRepMat}
Dado el siguiente sistema de ecuaciones lineales:
  \[
        \systeme[x_1 x_2 x_3]{x_1+5x_2-4x_3=3{,}, x_1-2x_2+x_3=1{,}, 3x_1+x_2-2x_3=-3.}
    \]
Su representación matricial es
\[
    Ax = b
\]
siendo 
 \[
        A=\begin{pmatrix}
        1 & 5 & -4\\
        1& -2 & 1\\
        3&1&-2\\
        \end{pmatrix}
\]
su matriz de coeficientes y 
\[
        b=\begin{pmatrix}
        3\\ 1 \\ -3
        \end{pmatrix}
        \texty
        x=\begin{pmatrix}
        x_1\\ x_2 \\ x_3
        \end{pmatrix}
\]
sus matrices columna de constantes y de incógnitas, respectivamente.
\end{ejem}

\begin{ejer}\hspace{0pt} 
Dada la representación matricial de un sistema de ecuaciones lineales:
\[
    Ax = b
\]
siendo
\[
    A = 
    \begin{pmatrix}
    1 & 5 & -4\\
    2& -2 & 3\\
    3&7&0\\
    0&1&1
    \end{pmatrix},
    \qquad
    x = 
    \begin{pmatrix}
    x_1\\ x_2\\ x_3
    \end{pmatrix}
    \texty
    b =
    \begin{pmatrix}
    1\\ 2 \\ -1\\3
    \end{pmatrix}.
\]
Escriba el sistema de ecuaciones lineales correspondiente.
\end{ejer}

\begin{defi}[Matriz ampliada]
    Sean $A\in\Mat[\K]{m}{n}$ y $B\in\Mat[\K]{m}{p}$, se define la matriz ampliada de $A$ y $B$ al elemento de $\Mat[\K]{m}{(n+p)}$ dado por:
    \[
        \begin{pmatrix}
        a_{11}&a_{12}&\cdots&a_{1n}&|&b_{11}&b_{12}&\cdots&b_{1p}\\
        a_{21}&a_{22}&\cdots&a_{2n}&|&b_{21}&b_{22}&\cdots&b_{2p}\\
        \vdots&\vdots&\ddots&\vdots&|&\vdots&\vdots&\ddots&\vdots\\
        a_{m1}&a_{m2}&\cdots&a_{mn}&|&b_{m1}&b_{m2}&\cdots&b_{mp}
        \end{pmatrix}.
    \]
    y se la denota por $(A|B)$.
\end{defi}

\begin{ejem}\hspace{0pt} 
Dadas las siguientes matrices:
\[
    A=\begin{pmatrix}
        1 & 3\\
        1& 2
\end{pmatrix}
\texty
B=\begin{pmatrix}
        1 & 0\\
        0& 1
\end{pmatrix}.
\]
Su matriz ampliada es:
\[
    (A|B)=
    \begin{pmatrix}
    1& 3&|&1 & 0\\
    1& 2&|&0&1
    \end{pmatrix}.
\]
\end{ejem}
\begin{defi}
    Dado el sistema de ecuaciones lineales en forma matricial
    \[
        Ax=b
    \]
    con $A\in \Mat[\K]{m}{n}$ y $b\in\Mat[\K]{m}{1}$, se dice que 
    
    \[
        (A|b)
    \]
    
    es la matriz ampliada asociada al sistema.
\end{defi}

\begin{ejem}
\hspace{0pt} 
Considere el sistema de ecuaciones lineales 
\[
        \systeme[x_1 x_2 x_3]{x_1+5x_2-4x_3=3{,}, x_1-2x_2+x_3=1{,}, 3x_1+x_2-2x_3=-3.}
\]
se tiene que su matriz ampliada es:
\[
    \begin{pmatrix}
        1&5&-4&|&3\\
        1&-2&1&|&1\\
        3&1&-2&|&-3
        \end{pmatrix}.
\]
\end{ejem}

\begin{ejer}\hspace{0pt} 
Dados los siguientes sistemas de ecuaciones lineales:
  \[
        \systeme[x_1 x_2 x_3]{x_1+5x_2-4x_3=3, x_1-2x_2+x_3=1, 3x_1+x_2-2x_3=-3}
        \texty
        \systeme[x_1 x_2 x_3]{x_1+x_2+x_3=0, x_1-6x_2+2x_3=-1.}
    \]
Determinar sus matrices asociadas.    
\end{ejer}


\subsection{Resolución de sistemas lineales}

\begin{defi}[Solución de un sistema]
Dados $A\in\Mat[\K]{m}{n}$, $b \in\Mat[\K]{m}{1}$ y $x \in\Mat[\K]{n}{1}$, considere la representación matricial $Ax = b$ de un sistema de ecuaciones lineales con $m$ ecuaciones y $n$ incógnitas. 
Se dice que una solución de este sistema es todo $z\in\Mat[\K]{n}{1}$ tal que $Az=b$.
\end{defi}

\begin{teo}
    Sean $A,C\in \Mat[\K]{m}{n}$ y $b,d\in \Mat[\K]{m}{1}$, se tiene que los sistemas de ecuaciones lineales
    \[
        Ax=b
        \texty
        Cx=d
    \]
    tienen las mismas soluciones si y solo si
    \[
        (A|b)\sim (C|d),
    \]
    es decir, si las matrices aumentadas de los sistemas son equivalentes por filas. 
\end{teo}

\noindent El procedimiento de reducción de Gauss-Jordan para resolver el sistema lineal $Ax = b$ es el siguiente.

\begin{description}
\item[Paso 1.] Formar la matriz aumentada $(A|b)$.
\item[Paso 2.] Transformar la matriz aumentada $(A|b)$ a su forma escalonada reducida por filas $(C|d)$ mediante operaciones elementales por fila.
\item[Paso 3.] Para cada fila distinta de cero de la matriz $(C|d)$, se despeja la incógnita correspondiente a la entrada principal de cada fila. Las filas que constan completamente de ceros se pueden ignorar, pues la ecuación correspondiente será satisfecha por cualquier valor de las incógnitas.
\end{description}

\noindent El procedimiento de eliminación de Gauss para resolver el sistema lineal $Ax = b$ es el siguiente.
\begin{description}
\item[Paso 1.] Formar la matriz aumentada $(A|b)$.
\item[Paso 2.] Transformar la matriz aumentada $(A|b)$ a su forma escalonada por filas $(C|d)$ mediante operaciones elementales por fila.
\item[Paso 3.] Resolver el sistema lineal correspondiente a $(C|d)$. Las filas constan completamente de ceros se pueden ignorar, pues la ecuación correspondiente será satisfecha por cualquier valor de las incógnitas.
\end{description}

\begin{ejem}\hspace{0pt} 
Los sistemas de ecuaciones lineales siguientes
\[
    \systeme[x_1  x_2]{2x_1 +x_2 = 2,
        x_1 -2x_2 = 1}
    \texty
    \systeme[x_1 x_2]{4x_1 +2x_2 = 4,
        7x_1 +x_2 = 7}
\]
tienen la misma solución, esta es $\begin{pmatrix}1\\0\end{pmatrix}$. En efecto, observemos que las matrices ampliadas de los sistemas son 
\[
    \begin{pmatrix}
    2&1&|&2\\
    1&-2&|&1
    \end{pmatrix}
    \texty
    \begin{pmatrix}
    4&2&|&4\\
    7&1&|&7
    \end{pmatrix}
\]
respectivamente, notemos también que
\begin{align*}
    \begin{pmatrix}
    2&1&|&2\\
    1&-2&|&1
    \end{pmatrix}
    &\sim \begin{pmatrix}
    2&1&|&2\\
    7&1&|&7
    \end{pmatrix}
    && 3F_{1} + F_{2} \to F_{2}\\
    &\sim \begin{pmatrix}
    4&2&|&4\\
    7&1&|&7
    \end{pmatrix}
    && 2F_{1} \to F_{1},
\end{align*}
es decir, las matrices ampliadas son equivalentes por filas, para obtener la solución del sistema utilizamos el procedimiento de reducción de Gauss-Jordan desde el paso 2, dado que ya obtuvimos a la matriz ampliada,
\begin{align*}
    \begin{pmatrix}
    2&1&|&2\\
    1&-2&|&1
    \end{pmatrix}
    &\sim \begin{pmatrix}
    1&-2&|&1\\
    2&1&|&2
    \end{pmatrix}
    && F_{1} \leftrightarrow F_{2}\\
    &\sim \begin{pmatrix}
    1&-2&|&1\\
    0&5&|&0
    \end{pmatrix}
    && -2F_{1} + F_{2} \to F_{2}\\
    &\sim \begin{pmatrix}
    1&-2&|&1\\
    0&1&|&0
    \end{pmatrix}
    && \frac{1}{5}F_{2} \to F_{2}\\
    &\sim \begin{pmatrix}
    1& 0&|&1\\
    0&  1&|&0
    \end{pmatrix}
    && 2F_{2} + F_{1} \to F_{1},
\end{align*}
lo que implica que la solución de los dos sistemas es $\begin{pmatrix}1\\0\end{pmatrix}$.
\end{ejem}

\begin{ejer}\hspace{0pt}
Considere a los sistemas de ecuaciones lineales siguientes:
\[
    \systeme[x_1 x_2]{x_1 +x_2 = 2,
        3x_1 -5x_2 = - 2}
    \texty
    \systeme[x_1 x_2]{2x_1 -x_2 = 1,
        x_1 -2x_2 = - 1.}
\]
Muestre que son sistemas equivalentes y obtenga la solución de estos.
\end{ejer}
\end{document}